% Hamilton's principle: stationary action among paths with fixed endpoints.
\begin{tikzpicture}[
  every node/.style={font=\small},
  vert/.style={fill=blue!75!black, circle, inner sep=2pt},
]
  % Endpoints
  \node[vert, label=below left:$q_0$]  (q0) at (-3.5, 0)    {};
  \node[vert, label=below right:$q_1$] (q1) at ( 3.5, 0)    {};

  % Physical trajectory (the critical path)
  \draw[red!75!black, very thick] (q0.center) .. controls (-1.5, 1.6) and (1.5, 1.6) .. (q1.center);
  \node[red!75!black] at (0, 1.85) {physical $\gamma$};

  % Variations (other admissible paths)
  \draw[gray!50!black, thin, dashed] (q0.center) .. controls (-1.5, 0.5) and (1.5, 0.5) .. (q1.center);
  \draw[gray!50!black, thin, dashed] (q0.center) .. controls (-1.5, 2.5) and (1.5, 2.5) .. (q1.center);
  \draw[gray!50!black, thin, dashed] (q0.center) .. controls (-1.5, 0.0) and (1.5, 1.2) .. (q1.center);

  \node[gray!50!black, font=\footnotesize] at (0, 2.7) {variations};
  \node[gray!50!black, font=\footnotesize] at (0, 0.2) {variations};

  % Action box
  \node[draw=red!50!black, fill=red!5, rounded corners=4pt, inner sep=4pt] at (0, 4.0)
    {$\delta S[\gamma] = 0 \;\Longleftrightarrow\; \frac{d}{dt}\frac{\partial L}{\partial \dot q} - \frac{\partial L}{\partial q} = 0$};

  \begin{scope}[on background layer]
    \fill[blue!4, rounded corners=12pt] (-4.5, -0.9) rectangle (4.5, 4.7);
  \end{scope}

  \node[align=center] at (0, -1.4)
    {\itshape Hamilton's principle: physical $\gamma$ is a critical point of $S = \int L\, dt$ among paths with fixed endpoints.};
\end{tikzpicture}
